\subsubsection{アーキテクチャ設計}
アーキテクチャ設計とは、設計原則と方法を使い、ソフトウェアアーキテクチャを作成し、記録する活動である。ここでは、システムの主要構成要素、その責任、性質、インタフェース、相互作用、環境との関係を決める。

典型的な関心事には、次のものがある。
\begin{itemize}
  \item 全体のアーキテクチャ様式と計算パラダイムである。
  \item システムを主要構成要素へ分割する方法である。
  \item 構成要素間の通信、同期、データ交換、制御の流れである。
  \item 関心事と設計責任をどの構成要素に割り当てるかである。
  \item 構成要素のインタフェースである。
  \item 性能、拡張性、資源消費、信頼性への対応である。
  \item 安全性やセキュリティのような全体横断的関心事への対応である。
\end{itemize}

アーキテクチャ設計は、\keyterm{分析}、\keyterm{統合}、\keyterm{評価}の三つの活動を反復する。
